\section{Experiments}
\label{sec:experiments}

We have three sequential experiments.
Section~\ref{sec:exp1} validates the boundary-regret estimator in synthetic settings with closed-form ground truth.
Section~\ref{sec:exp2} tests whether residual predictive expertise and residual decision value diverge on a real human--AI dataset across reward structures.
Section~\ref{sec:exp3} asks whether an ex ante boundary-regret score allocates scarce human review better than standard baselines under asymmetric rewards.

\subsection{Experiment 1: Estimator and Geometry Validation}
\label{sec:exp1}

\paragraph{Question.}
Does the plug-in estimator recover known boundary regret, and remain zero when a safe action dominates the reachable belief region?

\paragraph{Setup.}
We simulate finite-state Bayesian decision problems with known prior, reward matrix, and channels $P(m\mid y), P(h\mid y)$ ($M \perp H \mid Y$).
For each sampled $(y,m,h)$ we compute the exact posteriors $b_m=P(y\mid m)$ and $b_{m,h}=P(y\mid m,h)$, the model-only action $a_m$, and instance-level boundary regret
\[
  \BR(m,h) \;=\; V(b_{m,h}) - r_{a_m}\cdot b_{m,h}.
\]
The estimator does not see the true channels: it estimates $\hat P(m\mid y), \hat P(h\mid y), \hat p(y)$ from finite data via 5-fold cross-fitting with Laplace smoothing, then computes plug-in $\BRhat$ on held-out instances ($N=10{,}000$ per condition).

\paragraph{Conditions.}
The grid varies $G\in\{2,3,5\}$ decision groups and $K\in\{G,2G\}$ latent states, covering the minimal boundary case ($G=2$), the smallest genuinely multi-action case ($G=3$), and a moderate stress test ($G=5$); $K=2G$ checks that estimation does not degrade when multiple latent states share a decision group.
Three roles: C1 is a positive control (actions reward state groups, human signal crosses boundaries), C2 a weak null (prior deep inside one decision region, weak human channel), and C3 a strong null (highly predictive human signal but a safe action dominates every reachable posterior).
For C3 we set $\tau=(1+U)/2$, where $U$ is the maximum reachable posterior mass on any decision group under the exact channels; $\tau>U$ guarantees no reachable posterior can make a non-safe action optimal while the human signal retains its predictive accuracy.

\paragraph{Results.}
Table~\ref{tab:exp1-core} reports metrics averaged across the six configurations; rank correlation is reported only where $\BR$ varies (in both nulls $\BR\equiv 0$, so we report false-positive rate instead).

\begin{table}[htbp]
\centering
\caption{
    Synthetic validation of $\BRhat$ averaged over $G\in\{2,3,5\}$, $K\in\{G,2G\}$.
    C3's $\Delta\mathrm{LL}$ (0.595) exceeds C1's because the human signal is strongly predictive of group membership, yet no posterior crosses the certified threshold.
    $\Delta\mathrm{LL}$: log-loss improvement from $b_m$ to $b_{m,h}$.
}
\label{tab:exp1-core}
\small
\begin{tabular}{@{}llrrrrr@{}}
\toprule
Condition & Role & $\rho(\BRhat,\BR)$ & FPR & $\%\BR{>}0$ & $\%\BRhat{>}0$ & $\Delta$LL \\
\midrule
C1 & Positive control & 0.901        & 0.000 & 37.9 & 37.9 & 0.408 \\
C2 & Weak null        & \text{n/a}   & 0.000 &  0.0 &  0.0 & 0.019 \\
C3 & Strong null      & \text{n/a}   & 0.000 &  0.0 &  0.0 & \textbf{0.595} \\
\bottomrule
\end{tabular}
\end{table}

The estimator separates predictive value from decision value.
In C1, $\BRhat$ ranks true boundary regret well ($\rho=0.901$) across the 37.9\% of instances where crossings occur (per-configuration details in Appendix~\ref{app:exp1-detail}); FPR is zero in both nulls.
C3 is the critical case: the human signal achieves the highest log-loss improvement of any condition ($\Delta\mathrm{LL}=0.595$), yet $\BRhat$ assigns zero boundary regret to every instance because no posterior crosses the certified threshold.
High predictive accuracy on task-relevant features does not imply decision value when a safe action dominates the reachable belief region.

% ==============================================================

\subsection{Experiment 2: Residual Expertise Does Not Equal Decision Value}
\label{sec:exp2}

\paragraph{Question.}
Does the correlation between residual predictive expertise and boundary regret invert across reward regimes on the same model, readers, and data?

\paragraph{Setup.}
We pair a pretrained DenseNet-121 chest X-ray classifier \citep{Huang2017,Cohen2022} with five board-certified radiologists of the CheXpert reader study \citep{Irvin2019} across five official competition conditions (Atelectasis, Cardiomegaly, Consolidation, Edema, Pleural Effusion), yielding $N=25$ reader--task pairs on 500 frontal radiographs.
For each pair, the ground-truth label $y$ is the majority vote of the \emph{remaining four} readers, which eliminates structural $h$--$y$ contamination and ensures $y \perp h$.
Model beliefs $b_x$ are temperature-scaled per pair on a 30\% calibration split; evaluation uses the remaining 70\% (${\approx}350$ images per pair).
The human-augmented posterior $\bpost_{x,h} \approx P(Y=1 \mid x, h)$ is estimated via 5-fold cross-fitted logistic regression on features $[\operatorname{logit}(b_x),\, h,\, h \cdot \operatorname{logit}(b_x)]$.
Spearman correlations between residual predictive expertise (log-loss gain over $b_x$) and $\BRhat$ are reported across the 25 pairs under each reward regime; label protocol details in Appendix~\ref{app:exp2-protocol}.

\paragraph{Reward Regimes.}
We evaluate three reward matrices $r_{a,y}$, each placing the decision threshold $r_{0,0}/(r_{0,0}+r_{1,1})$~\citep{PaukerKassirer1980} in a distinct region of the model belief distribution.
R1 is the identity ($r_{0,0}=r_{1,1}=1$, errors $0$), threshold $1/2$, in the bulk.
R2 ($r_{0,0}=1$, $r_{1,1}=5$, errors $0$) penalises false negatives five times as heavily as false positives, threshold $1/6$, below the bulk.
R3 mirrors R2 ($r_{0,0}=5$, $r_{1,1}=1$), threshold $5/6$, in the sparse upper tail.
R3 is the critical negative control: Theorem~\ref{thm:facet-bridge} predicts that when the threshold sits in the tail, facet crossings become rare, boundary-regret mass collapses to near zero, and residual predictive expertise decouples from decision value regardless of how much the reader improves the model's log-loss.

\begin{table}[htbp]
\centering
\caption{Boundary-regret diagnostics on $N=25$ reader--task pairs: actionable BR mass collapses from R1 to R3 and BR concentrates near the active facet in all regimes.
  $n_{\mathrm{v}}$: pairs where $\BRhat\not\equiv 0$; $\bar\rho_{\mathrm{facet}}$: mean Spearman between facet distance and $\BRhat$.
}
\label{tab:exp2}
\small
\begin{tabular}{@{}llrrrrr@{}}
\toprule
Reward & Role & $\bar\rho_{\mathrm{LL}}$ & $\rho{>}0/n_{\mathrm{v}}$ & $\bar\rho_{\mathrm{facet}}$ & Top-decile zero-$\BRhat$ & $n_{\mathrm{v}}$ \\
\midrule
R1 (thr.\ $1/2$) & Primary result    & $\mathbf{+0.477}$ & $25/25$ & $-0.599$ & $11.2\%$          & $25$ \\
R2 (thr.\ $1/6$) & Intermediate      & $-0.051$          & $15/25$ & $-0.866$ & $28.3\%$          & $25$ \\
R3 (thr.\ $5/6$) & Negative control  & $-0.112$          & $0/7$   & $-0.221$ & $\mathbf{97.4\%}$ & $7$  \\
\bottomrule
\end{tabular}
\end{table}

\paragraph{Reward-Conditional Dissociation.}
The correlation between residual predictive expertise (log-loss gain over the model alone) and boundary regret inverts with the reward (Table~\ref{tab:exp2}): $\bar\rho_{\mathrm{LL}}=+0.477$ under R1 (25/25 pairs positive), $-0.051$ under R2 (15/25), and $-0.112$ under R3 (0/7; sign test $p=0.008$, with the remaining 18 of 25 pairs structurally $\BRhat\equiv 0$).
Same model, same readers, same data; only the reward changes.

The mechanism is geometric.
Under R1, threshold $1/2$ sits in a dense region of model beliefs, so a single reader update routinely crosses a facet and predictive expertise aligns with decision value.
Under R3, the false-positive facet sits in the tail and no posterior update reaches it: only 7 of 25 pairs carry any non-zero $\BRhat$, and $97.4\%$ of top-expertise-decile cases (90th percentile of model--reader disagreement $1-b_x[h]$) have $\BRhat=0$.
The R3 correlation $\bar\rho=-0.112$ on these 7 pairs is directionally consistent with the geometric prediction; the collapse of $n_\mathrm{v}$ from 25 to 7 is the direct confirmation that boundary-regret mass has evacuated.
R2 is intermediate ($28.3\%$ top-decile zero-$\BRhat$).
Facet distance correlates negatively with $\BRhat$ under every reward ($-0.599, -0.866, -0.221$ for R1, R2, R3): boundary regret concentrates near the active facet regardless of how predictive the reader is (Theorem~\ref{thm:facet-bridge}).
R2's stronger concentration follows from the same geometry: with the threshold at $1/6$ nearly all decisions default to treat, so the rare facet crossings come precisely from beliefs already near $1/6$. R1's centrally located threshold produces no such concentration.
We call this \emph{reward-conditional dissociation}: a consequence of reward geometry, not a calibration failure.

Appendix~\ref{app:exp2-robustness} shows the inversion and dissociation hold across in- and off-domain pretrained checkpoints (varying the belief geometry) and readers not involved in ground-truth labelling.
As a theory-confirming negative control, Appendix~\ref{app:exp2-cifar10h} evaluates CIFAR-10H: on highly accurate classifiers whose softmax probabilities rarely straddle a reward facet, $\bar\rho$ collapses to near-zero across all rewards ($+0.015$ under R1 vs.\ $+0.477$ on CheXpert R1), confirming that the dissociation signal depends on the model's beliefs straddling the active facet, exactly as Theorem~\ref{thm:facet-bridge} predicts.

\paragraph{Implication for Predictive Audits.}
A residual-expertise audit in the style of \citet{AlurEtAl2023,AlurRaghavanShah2024} would find substantial predictive value here (mean log-loss gain $0.387$ nats across all 25 pairs) and recommend deploying the readers.
Under R3, $97.4\%$ of those same cases carry zero decision value.
The \emph{estimand shift} from predictive to decisional complementarity is not a refinement; it changes the sign of the deployment recommendation.
If this regime-conditional inversion is theorem-predicted, an ex ante $\BRhat$ score should inherit it: dominating reward-aware geometric proximity where boundary-regret mass is reachable, converging with it where mass is absent.
Section~\ref{sec:exp3} tests this directly.

% ==============================================================

\subsection{Experiment 3: Review Allocation under Scarce Human Input}
\label{sec:exp3}

\paragraph{Question.}
Given a fixed human-review budget, can an ex ante boundary-regret score identify the cases where reader input will improve realized decision utility, rather than improve prediction or flag model uncertainty?
The prediction is regime-conditional: it should beat reward-aware but human-blind boundary proximity when reader updates can cross reward-relevant decision boundaries, and converge to it when such crossings are sparse or absent.


\paragraph{Setup.}
We use the same 25 reader--task pairs as in Section~\ref{sec:exp2}.
For each case, the model-only belief $b_x$ induces a model-only action $a_x$.
For each review budget $q\in\{5\%,\,10\%,\,20\%,\,50\%\}$, a policy ranks cases using only pre-review information.
Reviewed cases observe the reader response $h$ and take the human-updated action $\argmax_a r_a\cdot\bpost_{x,h}$; unreviewed cases retain $a_x$.
Utility gain is the held-out realized reward improvement over this no-review baseline.
All learned policy scores are 5-fold cross-fitted on the same fold structure used to estimate $\bpost_{x,h}$, so review priority is always computed out of fold.

\paragraph{Policies.}
The proposed policy, $\BRhat$, is an ex ante boundary-regret score: it predicts possible reader responses from $x$, predicts the corresponding human-updated posterior, and scores a case by the expected decision value of that update under the current reward.
We compare it against six baselines.
\textbf{Margin} is a reward-aware distance to the active decision facet, our primary comparison baseline: it captures reward geometry without estimating the human signal $P(h\mid x)$.
\textbf{Residual} is a residual-expertise score, the deployable analogue of a predictive audit \citep{AlurEtAl2023}.
\textbf{Entropy} is reward-blind model uncertainty.
These policies span two axes: Margin is reward-aware but human-blind, Residual is human-aware but reward-blind, Entropy is blind to both, and $\BRhat$ is both reward-aware and human-aware.
\textbf{L2D} is the learned full-coverage competitor: like $\BRhat$ it is both reward-aware and human-aware, but reaches that coverage end-to-end through a cross-fitted deferral classifier \citep{MozannarSontag2020,MaoMohri2023} with reward-sensitive labels $\mathbf{1}[R[a_{x,h},y]>R[a_x,y]]$, rather than through the reward's geometry and the reader-channel posterior.
\textbf{Random} is the uninformed baseline.
\textbf{Oracle} is the hindsight upper bound using realized labels.

\paragraph{Results.}

\begin{table}[t]
\centering
\caption{Section~\ref{sec:exp3}: mean utility gain over no-review baseline, $N\!=\!25$ reader--task pairs.
  W/L/T $=$ wins/losses/ties vs.\ Margin across the 25 pairs; decisive $=$ pairs where policies diverge.
  R3 is reported at $q\!=\!20\%$ only as a \emph{negative control}; the full R3 table is in Appendix~\ref{app:exp3-r3}.
  Under R1 (binary symmetric reward, threshold $1/2$), Margin and Entropy are rank-equivalent: both are strictly monotone in $|b_x - 1/2|$ for $K\!=\!2$, so they induce identical rankings and report identical gains.}
\label{tab:exp3}
\footnotesize
\begin{tabular}{@{}ll rrrrr r@{}}
\toprule
Reward & $q$ & $\BRhat$ & Margin & Entropy & Residual & L2D & W/L/T (sign-test $p$) \\
\midrule
R1 & $5\%$  & $\mathbf{+0.0473}$ & $+0.0395$ & $+0.0395$ & $+0.0226$ & $+0.0151$ & $20/1/4\ (p<0.0001)$ \\
R1 & $10\%$ & $\mathbf{+0.0893}$ & $+0.0743$ & $+0.0743$ & $+0.0465$ & $+0.0306$ & $24/1/0\ (p<0.0001)$ \\
R1 & $20\%$ & $\mathbf{+0.1634}$ & $+0.1438$ & $+0.1438$ & $+0.0917$ & $+0.0714$ & $22/1/2\ (p<0.0001)$ \\
R1 & $50\%$ & $\mathbf{+0.3337}$ & $+0.2797$ & $+0.2797$ & $+0.2545$ & $+0.2510$ & $25/0/0\ (p<0.0001)$ \\
\midrule
R2 & $5\%$  & $+0.0514$ & $+0.0490$ & $+0.0429$ & $+0.0317$ & $+0.0474$ & $3/0/22\ (p=0.25)$ \\
R2 & $10\%$ & $+0.0992$ & $+0.0965$ & $+0.0768$ & $+0.0621$ & $+0.0944$ & $3/1/21\ (p=0.63)$ \\
R2 & $20\%$ & $+0.1920$ & $+0.1889$ & $+0.1486$ & $+0.1270$ & $+0.1886$ & $3/1/21\ (p=0.63)$ \\
\midrule
R3$^\dagger$ & $20\%$ & $+0.0026$ & $+0.0029$ & --- & --- & --- & $0/2/23$ (neg.\ control) \\
\bottomrule
\end{tabular}
{\footnotesize $^\dagger$R3 negative control: $\BRhat$ and Margin are statistically indistinguishable ($p\geq 0.5$, W/L/T $=0/2/23$ at $q\!=\!20\%$); full table in Appendix~\ref{app:exp3-r3}.}
\end{table}

\paragraph{Regime-conditional inversion.}
Three reward matrices, one model, twenty-five reader--task pairs.
The same scoring rule produces a regime-conditional inversion of allocation advantage: $\BRhat$ dominates Margin under R1 at every budget (sign-test $p<0.0001$), is indistinguishable from Margin under R2 (sparse actionable mass at threshold $1/6$), and converges to Margin under R3 (only 7 of 25 pairs carry non-zero mass).
The pattern across rewards is the result; the wins are the evidence.

\paragraph{R1 (symmetric reward, threshold $1/2$).}
$\BRhat$ wins decisively at every budget (sign-test $p<0.0001$ throughout).
The four budget tallies $20/1/4$, $24/1/0$, $22/1/2$, $25/0/0$ yield decisive-win rates of $95\%$, $96\%$, $96\%$, $100\%$.
At $q=50\%$, $\BRhat$ wins all 25 pairs; at lower budgets, ties occur where both policies select the same case set, and decisive-win rates remain $95$--$96\%$.

\paragraph{R2 (FN-penalized reward, threshold $1/6$).}
The reward boundary at $1/6\approx 0.167$ sits closer to the belief mode than R3 but further than R1.
$\BRhat$ is indistinguishable from Margin at all tested budgets ($3/0/22$ to $3/1/21$; sign-test $p\geq 0.25$).
The near-null result reflects sparse actionable mass at threshold $1/6$: with LOO-corrected labels, the posterior updates that cross the R2 facet are too infrequent to produce a reliable allocation advantage over reward-aware Margin.

\paragraph{R3 (FP-penalized reward, threshold $5/6$).}
The reward boundary sits at $5/6\approx 0.833$, in the low-density tail of model beliefs.
Eighteen of twenty-five pairs have $\BRhat\equiv 0$ across the entire evaluation set (no reachable boundary crossing under R3's reward); among the 7 pairs with non-zero mass, $\BRhat$ is statistically indistinguishable from Margin at every budget ($p\geq 0.5$ throughout), with mean utility gains within $0.0005$ of Margin's at each budget ($q=20\%$: $\BRhat=+0.0026$ vs.\ Margin $=+0.0029$).
The geometry already isolates the only feasible boundary crossings; this is Theorem~\ref{thm:facet-bridge} confirmed, not contradicted.
The full R3 table is in Appendix~\ref{app:exp3-r3}.

\paragraph{Human channel versus geometric proximity.}
Reward-aware distance to the active facet is a strong baseline: it already identifies which cases are near the decision boundary.
The $95$--$100\%$ paired wins over Margin under R1 (sign-test $p<0.0001$ at every budget) show that $\BRhat$'s estimate of $P(h\mid x)$ and $P(Y\mid x,h)$ adds reliable signal above that geometric distance where actionable mass is dense.
Where mass is absent (R3) or sparse (R2), $\BRhat$ and Margin converge, both outcomes predicted by Theorem~\ref{thm:facet-bridge}.
Appendix~\ref{app:exp3-robustness} extends the allocation experiment to all six DenseNet checkpoints (240 pairs); R1 dominance holds across models, while off-domain checkpoints produce denser R2 actionable mass and stronger $\BRhat$ wins at smaller budgets.
